\documentclass[10pt]{article}

\usepackage[utf8]{inputenc}

\usepackage[T1]{fontenc}

\usepackage{amsmath}

\usepackage{amsfonts}

\usepackage{amssymb}

\usepackage[version=4]{mhchem}

\usepackage{stmaryrd}

\usepackage{graphicx}

\usepackage[export]{adjustbox}

\graphicspath{ {./images/} }

\usepackage{bbold}



\begin{document}

Аналог Х.н. справедлив в банаховых пространствах [4]. Существует такая константа $C(p, q), 0<p, q<\infty$, что для любых элементов $x_{k}$ из банахова пространства $E$



\[

\|\| \sum_{k=1}^{\infty} x_{k} r_{k}(t)\left\|_{E}\right\|_{L_{p}} \leqslant C(p, q)\|\| \sum_{k=1}^{\infty} x_{k} r_{k}(t)\left\|_{E}\right\|_{L_{q}} \cdot

\]



Одно из многочисленных приложений $\mathrm{X}$. н.: если



\[

\sum_{k=1}^{\infty} a_{k}^{2}+b_{k}^{2}<\infty

\]



то для почти всех наборов \textbackslash pm 1 функция



\[

\sum_{k=1}^{\infty} \pm\left(a_{k} \cos k t+b_{k} \sin k t\right)

\]



принадлежит всем $L_{p}, p<\infty$ (см. [5]).



Jum.: [1] Х и нч и н А., "Math. Z.», 1923, Bd 18, S. 109116; [2] K a r l i n S., "Trans. Amer. Máth. Soc.», 1949, v. 66, p. 44-64; [3] Г а по шा К ин В. Ф., «Успехи матем. наук», циональные ряды, пер. с англ., М., 1973; [5] 3 и г м у н д А., Тригонометрические ряды, пер. с англ., т. 1, М., 1965.



XIHЧИНА TEOPEMA 1) $\mathrm{X}$. T. з а ци и р а с п р е д е л е н и й: любое распределение вероя'гностей $P$ допускает (в сверточной полугруппе распределений вероятностей) факторизацию



\[

P=P_{1} * P_{2},

\]



где $P_{1}$ - распределение класса $I_{0}$ (см. Безгранично дeлимых распределений разложение), а $\boldsymbol{P}_{2}$ - распределение, к-рое либо является вырожденным, либо представимо в виде свертки конечного или счетного множества неразложимых распределений. Факторизация $(*)$, вообще говоря, не единственна.



Х. т. была доказана А. Я. Хинчиным [1] для распределений на прямой, позднее выяснилось [2], что она справедлива для распределений на значительно более общих группах. Известен (см. [3], [4] широкий клласс топологич. полугрупп, включающий сверточную полугруппу распределений на прямой, в к-рых справедливы факторизационные теоремы, аналогичные X. т. Лит.: [1] Х и н ч и н А. Я., "Бюлл. МГУ. Секц. А», 1937, т. 1, в. 1, c. 6-17; [2] П а р т а с а p a т и F. Р., Р а н г a



\includegraphics[max width=\textwidth, center]{2023_07_09_1eef3130fdf1638080cfg-1}

verw. Geb.), 1968, Bd 9, H. 3, S. 163-95; [4] D avid s on R., verw. Geb.", 1968, Bd 9, H. 3, S. $163-95 ;$ 4] D av i s on R.,

там же, 1968, Bd 10, H. 2, S. $120-72$.



\begin{enumerate}

  \setcounter{enumi}{1}

  \item Х. т. о ди оф а н т о вых при б лиж е ния х - см. Диофантовых приближений метрическая теория.

\end{enumerate}



ХОДЖА ГИПОТЕЗА - предположение о том, что для любого гладкого проективного многообразия $X$ над полем $\mathbb{C}$ комплексных чисел и для любого целого $p \geqslant 0 \mathbb{Q}$-пространство $H^{2 p}(X, \mathbb{Q}) \cap H^{p, p}$, где $H^{p, p}$ компонента типа $(p, p)$ в разложении Ходжа



\[

H^{2 p}(X, \mathbb{Q}) \bigotimes_{\mathbb{Q}} \mathbb{C}=\bigoplus_{r=0}^{2 p} H^{r, 2 p-r},

\]



порождается классами когомологий алгебраич. циклов коразмерности $p$ на $X$. Эта гинотеза была выдвинута У. Ходжем [1].



В случае $p=1$ Х.г. равносильна Лефшеца теореме о когомологиях типа $(1,1)$. X.г. доказана также для следующих классов многообразий: 1) $X$-гладкое 4-мерное уни лине й ч а тое много о б разие, т. е. такое, что существует рациональное отображение конечной степени $P^{1} \times Y \longrightarrow X$, где $Y-$ гладкое многообразие (см. [2]). Унилинейчатыми многообразиями являются, напр., унирациональные многообразия и 4-мерные полные пересечения с обильным антиканонич. классом (см. [3]). 2) $X$-гладкая гичерповерхность Ферма простой степени (см. [4], [5]). 3) $X-$ прстое 5-мерное абелево многообразие (см. [6]). 4) $X-$ простое $d$-мерное абелево многообразие, причем Fnd $(X) \bigotimes_{\mathbb{Z}} \mathbb{R}=\mathbb{R}^{l}, \frac{d}{l}$ - нечетное число, или End $(X) \otimes_{\mathbb{Z}} \mathbb{R}=\left[M_{2}(\mathbb{R})\right]^{l}, \frac{d}{2 l}$ - нечетное чіс ло.



Jum.: [1] H o d g e W. V. D., "Proc. Intern. Congr. Math.» (Camb., 1950), 1952, v. 1, p. 182-92; [2] C o n t a A., M u rr e J. P., "Math. Ann.", 1978, Bd 238, S. 79-88: [3] и х ж е, «J. de géometrie algébrique d'Angers» (juillet, 1979), 1980, p. $129-$ 41 ; [4] R a n Z., "Comp. math.", 1980/81, v. 42, N N 1, p. $121-42$; p. 111-14; [6] Т а н к е е в С. Г., «Изв. АН СССР. Сеp. матем.», 1981, т. $45, N_{2} 4$ с. $793-823$. С. Г. Танкеев.



ХОДЖА МНОГООБРАЗИЕ - комплексное многообразие, на к-ром можно задать м е т р и к у X о д ж а, т. е. Кәлера метрику, фундаментальная форма к-рой определяет целочисленный класс когомологий. Компактное комплексное многообразие является X. м. тогда и только тогда, когда оно изоморфно гладкому алгебраич. подмногообразию комплексного проективного пространства нек-рой размерности (т е о р е м а Код а и ры о проективном в ло же ни и).



См. также Кәлерово многообразие.



Лum.: [1] Г р и фф и т с Ф., Х а р р и с Д ж., Принципы алгебраической геометрии, пер. с англ., т. 1, М., 1982.



ХОДЖА СТРУКТУРА в е с а $n$ (ч и с т а я)-Объент, состоящий из репетки $H_{\mathbb{Z}}$ в действительном векторном пространстве $H_{\mathbb{R}} \stackrel{=}{=} H_{\mathbb{Z}} \otimes \mathbb{R} \quad$ и разложения $H_{\mathbb{C}}=\bigoplus_{p+q=n} H^{p, q}$ комплексного векторного шространства $H_{\mathbb{C}}=H_{\mathbb{Z}} \otimes \mathbb{C}($ разложения Холжа). Прі этом должно выполняться условие $\bar{H}^{p, q}=H^{q, p}$, где черта означает комплексное сопряжение в $H_{\mathbb{C}}=$ $=H_{\mathbb{R}} \otimes_{\mathbb{R}} \mathbb{C}$. Другое описание разложения Ходжа состоит в задании убывающей фильтрации (ф и л ь т р ации и $\mathrm{X} о$ дж а) $\quad F^{r}=\bigoplus_{p} \geqslant r^{H^{p, q}}$ в $H_{\mathbb{C}}$ такой, что $\bar{F}^{s} \cap F^{r}=0$ при $r+s \neq n$. Тогда подпространства $H^{p, q}$ восстанавливаются по формуле $H^{p, q}=F^{p} \cap \bar{F}^{q}$.



Примером является Х.с. в пространство $n$-мерных когомологий $H^{n}(X, \mathbb{C})$ компактного кәлерова многообразия $X$, впервые изученная У. Ходжем (см. [1]). В этом случае подпространства $H^{p, q}$ описываются как пространства гармонических борм типа $(p, q)$ или как когомологии $H^{q}\left(X, \Omega^{p}\right)$ пучков $\Omega^{p}$ голоморфых дифференциальных форм [2]. Фильтрация Ходжа в $H^{n}(X, \mathbb{C})$ возникает из фильтрации комплекса пучков $\Omega=\sum_{p \geqslant 0} \Omega^{p}, n$-мерные гиперкогомологии к-рого изоморфны $H^{n}(X, \mathbb{C})$, подкомплексами $\sum_{r \geqslant p} \Omega^{r}$.



Более общим понятием является с м е ші а н н а я X.c. Это-объект, состоящий из решетки $H_{\mathbb{Z}}{ }^{\text {в }}$ $H_{\mathbb{R}}=H_{\mathbb{Z}} \otimes \mathbb{R}$, возрастающей фильтрации (фильтрации весов) $W_{n}$ в $H_{\mathbb{Q}}=H_{\mathbb{Z}} \otimes \mathbb{Q}$ и убывающей фильтрации (фильтрации Ходжа) $F^{p}$ в $H_{\mathbb{C}}=H_{\mathbb{Z}} \bigotimes \mathbb{C}$ таких, что на пространстве $\left(W_{n} / W_{n+1}\right) \otimes \mathbb{C}$ фильтрации $F^{p}$ и $\bar{F}^{p}$ определяют чистую X.c. веса $n$. Рассмотрена [3] смешанная X.c. в когомологиях комплексного алгебраич. многообразия (не обязательно компактного или гладкого), гак аналог структуры модуля Галуа в этальных когомологиях. Х.с. имеют важные приложения в алгебраич. геометрии (см. Отображение периодов) и в теории особенностей гладких отображений (см. [4]).



Jum.: [1] H o d g e W. V. D., The theorie and applications of harmonic integrals, 2 ed., Camb., 1952; [2] $\Gamma$ p и ф ф ф и т с Ф., $\mathbf{X}$ a p p и с Д ж., Принципы алгебраической геометрии, пер. Math.» (Vancouver, 1974), 1975, v. 1, p. 79 -85; [4] B а р ч е нк о А. Н., в сб.: Современные проблемы математики, т. 22 , М., 1983 , с. $130-66$ (Итоги науки и техники). См. также лит.к статье отобрансени периодов.



А. И. Овсеевич.





\end{document}